% sec-04: the Phase I technical objectives.  Table 19.
\section{The Phase I Technical Objectives}

\subsection{What the Pitch proposes}

The Pitch proposes a governed, advisory-only language model layer for robotic
oncologic surgery, and a verification method that makes its behavior checkable.
The layer runs on premises at a clinical site. It produces text recommendations
to the operating surgeon and nothing else: it holds no write credential to the
data capture system, no route to the robot control network, and no ability to
act on equipment \cite{onpremwhippl,phase1ind}. Every output is checked by two
independent models from different developers before it is shown, a practice the
company already follows in its own work \cite{tripartisan2026}, and every check
is written to an audit trail that can be replayed.

The innovation is the combination of three things that clinical AI systems
usually treat separately: a hard architectural boundary, enforced by network
segmentation and credential design rather than by instructions to the model;
cross-model verification, in which disagreement between independent models is a
signal to withhold advice rather than noise; and a credibility assessment framed
on the ASME V\&V 40 approach to computational models \cite{asmevv40}, extended
from simulation to advisory text.

\subsection{Five measures, each with a threshold}

Table 19 states the Phase I as five measures. The first four are the technical
risks the Pitch names; the fifth restates the physical limits of the deposited
protocol that the layer must never touch. Each row has a threshold a reviewer
can check at the end of the award.

\begin{apptable}
\begin{tabularx}{\textwidth}{>{\raggedright\arraybackslash}p{2.9cm} >{\raggedright\arraybackslash}p{5.0cm} >{\raggedright\arraybackslash}p{4.6cm} >{\raggedright\arraybackslash}X}
\headrow \thead{Objective} & \thead{Measure} & \thead{Threshold} & \thead{Source}\\
\midrule
Boundary integrity & Open paths across the boundary, tested on every build & Zero write credentials; no route to robot control & Protocol, IND \\
Verification yield & How often verifiers disagree, and flag a seeded error & A withholding threshold set from measured rates & Company method \\
Credibility & An ASME V\&V 40 battery, extended to advisory text & Twin at 81.9 over 55 tests; extension with limits & Deposited twin \\
Timing & Latency against each safety stop, on a hardware bench & No effect on 3 ms arm or 500 ms system stops & Protocol \\
Physical limits & Any advisory path to a force or vessel control & Zero paths to the 3 N, 18 N and five-vessel limits & Protocol \\
\end{tabularx}
\end{apptable}
\vspace{-0.60cm}
\tabcap{Table 19. The Phase I as five measures, each with a threshold a reviewer can check at\\
the end of the award and the deposited source it comes from; no patient data is used.}

\subsection{What the Phase I would produce}

A Phase I would produce three things: a verification method, the test suite that
exercises it on every build, and a public report of measured rates with their
limitations. The main technical risks are building synthetic cases realistic
enough to expose failure, choosing verifier models that fail independently
rather than together, and measuring rare events without patient data. Each is a
question of engineering and measurement, which is the kind of question NSF
funds.

\subsection{What is not proposed}

No clinical trial is proposed, and no patient data is used, identifiable or
de-identified. The deployed clinical model stays on premises at the site, and
any shared computing resource would see only synthetic cases and the company's
own code. The robotic Phase 1 the layer would one day sit inside remains an
open-label, single-arm, 3+3 dose escalation of perioperative daraxonrasib at
160, 220 and 300 mg, with a 28-day dose-limiting toxicity window, in up to 18
treated participants undergoing a staged eight-arm robotic
pancreaticoduodenectomy with 56 degrees of freedom and 640 sensor channels
\cite{onpremwhippl}. That trial belongs to a site investigator and to NIH and
private funding, as \S3 sets out.
